% Part: first-order-logic
% Chapter: tableaux
% Section: derivations

\documentclass[../../../include/open-logic-section]{subfiles}

\begin{document}

\iftag{FOL}
      {\olfileid[hi]{fol}{tab}{der}}
      {\olfileid[hi]{pl}{tab}{der}}

\olsection{\usetoken{P}{tableau}}

\begin{explain}
हम बता चुके हैं कि मान्यता क्या होती है और अनुमान-नियम दे चुके हैं।
इन्हीं से !!^{tableau}s आगमनात्मक रूप से बनते हैं: प्रत्येक !!{tableau} या
तो एक अथवा अधिक मान्यताओं से बनी अकेली शाखा होता है, या फिर किसी शाखा पर
अनुमान-नियमों में से एक लगाने पर किसी !!a{tableau} से प्राप्त होता है।
\end{explain}

\begin{defn}[!!^{tableau}]
मान्यताओं $\sFmla{S_1}{!A_1}$, \dots,
$\sFmla{S_n}{!A_n}$ (जहाँ प्रत्येक $S_i$ या तो $\True$ है या~$\False$) के
लिए !!^a{tableau}, !!{signed formula}s का ऐसा परिमित वृक्ष है जो निम्न
शर्तें पूरी करता है:
\begin{enumerate}
\item वृक्ष के सबसे ऊपर के $n$~!!{signed formula}s
  $\sFmla{S_i}{!A_i}$ हैं, एक के नीचे दूसरा।
\item वृक्ष का प्रत्येक ऐसा !!{signed formula} जो मान्यताओं में से नहीं है,
  अपने ऊपर की शाखा के किसी !!a{signed formula} पर अनुमान-नियम के सही प्रयोग
  से प्राप्त होता है।
\end{enumerate}
!!a{tableau} की शाखा \emph{बंद} है तभी और केवल तभी जब उसमें
$\sFmla{\True}{!A}$ और~$\sFmla{\False}{!A}$ दोनों हों; अन्यथा वह
\emph{खुली} है। ऐसा !!^a{tableau} जिसकी प्रत्येक शाखा बंद है, अपनी
मान्यताओं के समुच्चय के लिए \emph{बंद !!{tableau}} है। यदि कोई
!!a{tableau} बंद नहीं है, अर्थात् उसमें कम-से-कम एक खुली शाखा है, तो वह
\emph{खुला} है।
\end{defn}

\begin{ex}
  मान्यताओं का प्रत्येक समुच्चय अपने-आप में !!a{tableau} है, पर सामान्यतः
  वह बंद नहीं होगा। (स्पष्टतः वह तभी बंद है जब मान्यताओं में पहले से
  !!{signed formula}s का कोई युग्म $\sFmla{\True}{!A}$ और
  ~$\sFmla{\False}{!A}$ हो।)

  किसी !!a{tableau} (खुले या बंद) में उपस्थित किसी
  !!a{signed formula}~$!A$ पर अनुमान-नियमों में से एक लगाकर हम उससे नया,
  बड़ा सत्य-वृक्ष प्राप्त कर सकते हैं। नियम~$!A$ की उस उपस्थिति वाली हर
  शाखा के अंत में एक या अधिक !!{signed formula}s जोड़ देगा जिस पर हमने
  नियम लगाया है।

  उदाहरणार्थ, मान्यता $\sFmla{\True}{!A \land \lnot
    !A}$ पर विचार करें।
  केवल इस मान्यता से बना (खुला) !!{tableau} यह है:
  \begin{center}
    \begin{tableau}{}
      [\sFmla{\True}{\formula{A} \land \lnot \formula{A}}, just=\TAss]
    \end{tableau}{}
  \end{center}
  मान्यता पर $\TRule{\True}{\land}$ नियम लगाकर इससे नया !!{tableau}
  मिलता है। यह नियम हमें !!{tableau} में दो नई पंक्तियाँ,
  $\sFmla{\True}{!A}$ और $\sFmla{\True}{\lnot !A}$, जोड़ने देता है:
  \begin{center}
    \begin{tableau}{}
      [\sFmla{\True}{\formula{A} \land \lnot \formula{A}}, just=\TAss
        [\sFmla{\True}{\formula{A}}, just={\TRule{\True}{\land}[1]},
          [\sFmla{\True}{\lnot \formula{A}}, just={\TRule{\True}{\land}[1]}
          ]
        ]
      ]
    \end{tableau}{}
  \end{center}
  !!{tableau}s लिखते समय हम लगाए गए नियम दाईं ओर दर्ज करते हैं (जैसे,
  $\TRule{\True}{\land} 1$ का अर्थ है कि उस पंक्ति का !!{signed formula},
  $\TRule{\True}{\land}$ नियम को पंक्ति~$1$ के !!{signed formula} पर लगाने
  से प्राप्त हुआ है)। इस नए !!{tableau} में अब अतिरिक्त !!{signed formula}s
  हैं, किंतु नियम केवल एक पर लगाया जा सकता है
  ($\sFmla{\True}{\lnot !A}$ पर; यहाँ $\TRule{\True}{\lnot}$ नियम)। इससे
  निम्न बंद !!{tableau} मिलता है:
  \begin{center}
    \begin{tableau}{}
      [\sFmla{\True}{\formula{A} \land \lnot \formula{A}}, just=\TAss
        [\sFmla{\True}{\formula{A}}, just={\TRule{\True}{\land}[1]}
          [\sFmla{\True}{\lnot \formula{A}}, just={\TRule{\True}{\land}[1]}
            [\sFmla{\False}{\formula{A}}, just={\TRule{\True}{\lnot}[3]}, close]
          ]
        ]
      ]
    \end{tableau}{}
  \end{center}
\end{ex}

\end{document}
